% ============================================================================
%  05_hardware.tex — Hardware fusionado y componentes nuevos.
% ============================================================================
\section{Hardware}
\label{sec:hardware}

El montaje fusiona en un único ESP32 DevKit~C~v4 los componentes de Act1 y
Act2 y añade los dos elementos nuevos exigidos por el autodiagnóstico
(tabla~\ref{tab:componentes}). El cableado completo se describe en
\codein{diagram.json} y se reproduce en la figura~\ref{fig:wokwi}.

\begin{table}[H]
  \centering
  \small
  \begin{tabular}{p{0.27\textwidth} p{0.38\textwidth} p{0.27\textwidth}}
    \toprule
    \textbf{Componente} & \textbf{Función} & \textbf{Pines ESP32} \\
    \midrule
    \multicolumn{3}{l}{\emph{Sensores compartidos / heredados}}\\
    DHT22 \textbf{principal} & Temperatura y humedad (sensor primario). & 4 \\
    LDR                  & Iluminación (ADC1, \emph{input-only}).         & 34 \\
    HC-SR04              & Altura por ultrasonidos (clasificador).        & 5 (TRIG), 18 (ECHO) \\
    LCD 20$\times$4      & HMI por I\textsuperscript{2}C (0x27).          & 21 (SDA), 22 (SCL) \\
    RTC DS1307           & Reloj para \emph{timestamps} (I\textsuperscript{2}C 0x68). & 21 (SDA), 22 (SCL) \\
    PIR HC-SR501         & Presencia en cabina.                           & 23 \\
    Receptor IR TSOP     & Mando a distancia NEC.                         & 15 \\
    Botones P0\,--\,P3   & Llamada por planta (P4 sólo por IR; ver abajo).& 32, 33, 35, 16 \\
    Botón RESET          & Reset de contadores del clasificador.          & 14 \\
    Servo cabina         & Movimiento entre plantas (PWM LEDC).           & 19 \\
    Servo compuerta      & Desvío de cajas (PWM LEDC).                    & 13 \\
    LED rojo/azul/amar.  & Calentador, enfriador, lámpara PWM.            & 2, 12, 0 \\
    LED verde/rojo       & Ascensor / Escalera (clasificación).           & 25, 26 \\
    Buzzer pasivo        & Pitido de clasificación y alarma.              & 27 \\
    \midrule
    \multicolumn{3}{l}{\emph{Nuevos para el autodiagnóstico (Act3)}}\\
    DHT22 \textbf{reserva} & Sensor T/H redundante (failover automático). & 17 \\
    Potenciómetro          & Sensor de corriente de actuadores (simulado).& 39 (V\textsubscript{N}) \\
    \bottomrule
  \end{tabular}
  \caption{Componentes y asignación de pines del sistema integrado.}
  \label{tab:componentes}
\end{table}

\begin{figure}[H]
  \centering
  \IfFileExists{img/wokwi_diagram.png}{%
    \includegraphics[width=0.85\textwidth]{wokwi_diagram.png}%
  }{%
    \fbox{\textit{[Pendiente: captura de pantalla del montaje Wokwi]}}%
  }
  \caption{Montaje completo en Wokwi: clasificador, ascensor, HMI y la pareja
    de sensores nuevos del autodiagnóstico (DHT22 de reserva y potenciómetro).}
  \label{fig:wokwi}
\end{figure}

\subsection{Pines nuevos del autodiagnóstico}

\paragraph{DHT22 de reserva — GPIO17.}
Es la \emph{redundancia} de temperatura y humedad pedida por la sección de
autodiagnóstico del enunciado. GPIO\,17 se ha liberado retirando el pulsador
físico de la planta~4, ya que era el único GPIO digital bidireccional
disponible tras Act1 y Act2: el ESP32 sólo deja libres GPIO\,36 y~39, ambos
\emph{input-only}, no aptos para el bus bidireccional del DHT22.
\textbf{La planta~4 sigue siendo plenamente accesible} mediante el botón~4
del mando IR; los pulsadores físicos y el mando IR fueron siempre redundantes
entre sí, así que la operativa del ascensor no pierde ninguna planta.

\paragraph{Sensor de corriente — GPIO39 (V\textsubscript{N}).}
Un potenciómetro modela la \emph{salida ya acondicionada} de un sensor de
corriente real (resistencia \emph{shunt} más amplificador): en Wokwi no se
puede variar el consumo real de un LED, así que el potenciómetro actúa como
inyector de fallo de \emph{sobreconsumo} para la demostración, y al mismo
tiempo deja en el diagrama el divisor analógico que tendría su contrapartida
hardware en un montaje físico. GPIO\,39 (etiquetado V\textsubscript{N} en la
DevKit) es \emph{input-only} y pertenece a ADC1, evitando cualquier conflicto
con el LDR (GPIO\,34) o con Wi-Fi (que reservaría ADC2). El escalado lineal
es \SIrange{0}{800}{\milli\ampere} con umbral de sobreconsumo en
\SI{500}{\milli\ampere}.

Los pines de Act1 y Act2 se eligieron con la regla de \emph{conjuntos
disjuntos} salvo los sensores compartidos, lo que permitió la fusión sin
cambios de cableado. Los \emph{strapping pins} del ESP32
\cite{espressif_esp32} mantienen niveles correctos al arranque: los LEDs de
control parten apagados y el receptor IR queda en \codein{HIGH}, compatible
con el reposo del TSOP.
